\section{Remaining questions and the boundary of the results}
\label{sec:open}

The external Nishibuchi conjecture, in its stated parameter range, is resolved
by the arithmetic results above. The following questions concern extensions,
sharper classifications, and the original native seed-calculus program. They
are not missing steps in that resolved conjecture. The local analyticity and
differentiability theorem is also complete under its stated invertibility
hypothesis.

\subsection{Singular persistence beyond an invertible solved block}
\label{subsec:singular-open}

The weighted lifting theorem allows the solved Jacobian to be noninvertible
modulo $p$ provided that it is invertible over $\Qp$ with a controlled inverse.
This includes a substantial class of lifts often called singular modulo $p$.
The remaining problem concerns a different situation: the Jacobian is
rank-deficient even over $\Qp$, and the current inverse construction cannot
apply.

For a fixed polynomial system $F\in\Z[X_1,\ldots,X_d]^e$, let
\[
\mathcal L_{F,L}(\bar x)
=\{x'\in(\Z/p^{L+1}\Z)^d:
F(x')=0,\ x'\equiv\bar x\pmod{p^L}\}.
\]
The finite obstruction theorem already decides whether this set is empty
from the actual residual digit and reduced Jacobian. If it is nonempty, its
cardinality and affine kernel structure are determined. A vanishing class
at one level says only that this one step is possible.

For example, $X^2-5$ has the root $1$ modulo $2$, and both $1$ and $3$ lift it
to roots modulo $4$. Every odd square is $1$ modulo $8$, so neither branch
continues to a root modulo $8$. This direct arithmetic example illustrates
why a successful first correction is not a persistence theorem.

The open task is to find useful structural hypotheses and higher obstruction
data that control an entire singular branch, including its behavior under
coefficient and parameter deformation. One must establish those conditions
at successive levels, or establish a suitable all-level existence statement
from which compatible branches follow. The finite obstruction calculus gives
a way to test each level. It does not supply an unconditional singular
persistence result. Multiple branches, later obstruction, and changes in the
rank of the Jacobian must remain possible in any proposed general statement.

\subsection{Native 13.2: what presentation invariance must mean}
\label{subsec:native132}

The original conjecture asks for a canonical finite obstruction packet whose
vanishing is necessary and sufficient for a one-level lift in a finitely
presented arithmetic seed system. For a fixed finite polynomial presentation,
the paper now supplies the precise obstruction
\[
o_{F,\bar x}=[r]\in
\coker\!\left(J_{F,\bar x}:\F_p^d\longrightarrow\F_p^e\right),
\]
where $r$ is the residual digit. Its independence from the chosen integer
representative is proved. That independence is different from invariance
under a change of equations or presentation.

The distinction cannot be removed by declaring the raw cokernel canonical.
The presentations $F=(X)$ and $G=(X,0)$ describe the same zero locus over every
coefficient ring. At zero their Jacobians are, respectively,
\[
\F_p\xrightarrow{\,1\,}\F_p,
\qquad
\F_p\xrightarrow{\,(1,0)\,}\F_p^2.
\]
Their raw cokernels are $0$ and $\F_p$. Both distinguished obstruction classes
vanish, but the two ambient obstruction groups are not isomorphic. This is
an elementary presentation counterexample, not a refutation of every possible
intrinsic obstruction construction.

A complete version of Native 13.2 therefore needs the following data before
its full statement can be audited:
\begin{enumerate}
\item a category of seed objects, including how the operators, relations,
projectors, weights, and topology enter an object;
\item specified changes of presentation and a proof that they preserve the
actual lift problem;
\item an obstruction object that respects dependencies among the relations,
together with comparison maps or a universal property;
\item a precise meaning of ``finite'' and ``canonical'', and a proof that
vanishing is equivalent to the intended one-step continuation.
\end{enumerate}
For commutative algebra presentations, retaining the relation module rather
than all listed equations is a natural route. Extending such a construction
to the entire seed vocabulary requires an encoding and a proof that lifts
are preserved. The consolidation does not claim that missing implication.

\subsection{Native 13.3: finite compatibility, completion, and descent}
\label{subsec:native133}

The original proposed form is
\[
X\simeq X_2^{\mathrm{wt}}\times_C\prod_{p>2}'X_p.
\]
It leaves $X$, the distinguished subsets defining the restricted product,
the maps to $C$, the ambient category, and the universal meaning of
``smallest'' unspecified. In particular, one must decide whether $X$ is a
global object or a completed object.

There is now an exact linear compatibility result. If two local observations
$J:U\to V$ and $K:U\to W$ are linear maps from the \emph{same} correction
module, define
\[
C_{J,K}=U/(\ker J+\ker K).
\]
The attainable observations $\im J$ and $\im K$ have canonical labels in this
quotient. Equality of labels is equivalent to simultaneous realization by a
single correction. The joint image is exactly the fiber product over
$C_{J,K}$, with the proved linear universal property.
The declarations are
\begingroup\small\raggedright
\lean{Komorebi.SharedCorrection.labels_eq_iff_common},
\lean{Komorebi.SharedCorrection.joint_range_eq_fiber}, and
\lean{Komorebi.SharedCorrection.common_universal}.
\par\endgroup
This result assumes a common correction module and actual observation maps.
An arithmetic application must construct them and prove that their notion
of compatibility is the one required by the application.

\paragraph{Why the global/completed distinction is necessary.}
Consider constant sequences with integer initial value. Their global parameter
space is $\Z$, while their local parameter spaces are $\Z_p$. Under the
canonical local maps, the image of the global objects is the diagonal
\[
D=\{(n)_p:n\in\Z\}\subseteq\prod_p\Z_p.
\]
This diagonal is dense: a basic open set specifies finitely many congruences,
and the Chinese remainder theorem provides an integer satisfying them. It is
proper: the tuple with $2$-adic coordinate zero and all odd-adic coordinates
one is not the image of an integer. Indeed an integer zero in $\Z_2$ must be
zero, contradicting its $3$-adic coordinate.

Consequently this canonical global diagonal cannot be an equality locus
\[
\{(a,b):f(a)=g(b)\}
\subseteq \Z_2\times\prod_{p>2}\Z_p
\]
for continuous maps to a Hausdorff space $C$. Such an equality locus is closed,
as the inverse image of the closed diagonal in $C\times C$, whereas $D$ is
proper and dense. This elementary topological argument is retained from the
native progress note as a prose argument. It is not a new Lean-certified
result, and it concerns exactly the stated topological interpretation.

Finite CRT reconstruction and the familiar completed object
$\widehat\Z=\prod_p\Z_p$ address a different question. They do not recover
which local tuples arise from a single global integer. A complete form of
Native 13.3 must specify finite-level compatibility, transition maps and
completion, and any additional global descent condition. Recurrence,
spectral, or Frobenius information must appear as actual maps or constraints,
not only as a proposed description of $C$.

\subsection{Arithmetic refinements}
\label{subsec:arithmetic-open}

Write $U_{N,j}(n)=j!C_{N,j}(n)$ and $U_j=U_{4,j}$, and use
$v_p(0)=\infty$ when discussing valuation bounds. The following classification
questions remain open in this program.

\begin{enumerate}
\item \textbf{Higher dyadic differences.}
Determine the exact value of
\[
v_2\!\left(\Delta_u^kU_j(n)\right),\qquad k\geq2,
\]
with explicit treatment of zero differences. The all-order lower bounds and
the exact first-displacement theorem do not determine this value: subtraction
of two first differences can introduce further cancellation. The first
concrete target is order $2$ in one fixed address.

\item \textbf{Least periods beyond $N=4$.}
For further packet families, determine the least positive $P$ for which
$U_{N,j}(n+P)\equiv U_{N,j}(n)\pmod{p^r}$ for every $n$ in the claimed range.
The general sufficient periods are proved. Their minimality does not follow
from those sufficient bounds, and a period in packet index is a degree shift
of $NP$.

\item \textbf{Exact affine Mahler-kernel valuations.}
For the kernel of Section~\ref{subsec:polyphase-exposition}, determine exact
values of
\[
v_p\bigl(\Theta_{M,s}(r,k)\bigr),\qquad
\Theta_{M,s}(r,k)=[z^r](1+z)^s((1+z)^M-1)^k,
\]
beyond the explicit special families. The minimum valuation of the summands
gives a lower bound. Equality requires a separate argument excluding
cancellation.

\item \textbf{Exact classical dyadic tail values.}
The uniform truncation exponent $4M+3$ and its optimality are proved.
Individual term weights are also known. The valuation of the whole omitted
tail at an arbitrary even row may be higher, so an exact formula for all rows
requires more than taking the minimum of those term weights. Empty tails
must remain exactly zero.
\end{enumerate}

These questions refine existing estimates and exact cases. They do not reopen
the external conjecture or the sharp uniform classical dyadic bound.

\subsection{Historical material and interface work}
\label{subsec:historical-scope}

The original status-corrected Paper IX survives as a 28-page PDF. Its recovered
formalization supports the framework, Euler coordinates, Mahler differences,
unit coordinates, recurrence packets, and CRT results retained in this
manuscript. The original TeX, original figure-generation files, and complete
historical experiment package were not recovered with that PDF. Their older
figure and test counts are therefore historical records, not newly reproduced
computations in this consolidation.

A separate reconstructed Annals text survives as a retrieved textual copy.
It includes finite-field mode and subgroup inversion, Fourier invariance,
Carmichael and matrix transfer, Amice pullbacks, depleted analytic Euler
branches, and extensive polyphase formulas. The sharp classical dyadic kernel
and truncation claims from that text have since been connected to checked
Lean proofs. The finite polyphase formulas used for explanation here have
short explicit algebraic derivations. The remainder of that larger historical
text has not acquired blanket formal coverage merely by being preserved.
The source archive retains the text with its provenance, and the statement
map distinguishes these historical claims from the certified claims in the
present manuscript.

Remaining recovery work includes the precise older odd-prime edge identities,
their stated ranges, and the binary digit-sum form of the checked weights.
Recovering a formula, checking its prose proof, and locating an exact Lean
declaration are separate steps. Missing historical sources should remain
identified instead of being reconstructed from a remembered result.

There is also a narrower interface task. The new regularity theorem applies to
arbitrary ambient $\Qp$ polynomial coefficients and identifies the analytic
germ of every continuous ambient-field root selection under its hypotheses.
Several older rectangular charts were defined using integer coefficients,
integer centers, or subtype-valued maps. Comparing each such historical
interface explicitly with the new ambient formulation remains to be done.
It does not invalidate the proved analytic implicit-function theorem, and it
is not counted as a general singular-persistence result.
